\documentclass{article}
\usepackage[utf8]{inputenc}
\usepackage[T1]{fontenc}
\usepackage[margin=2cm]{geometry}
\usepackage{hyperref}
\sloppy
\begin{document}
\section*{Historical source}
Encyclopaedia Britannica scan, source row 268. Parse 3.5 extracted the text. Super 3.5 VL corrected it from an enlarged crop of each block. It can still contain recognition errors.
Page headers, footers, and page numbers are removed.

\subsection*{Untitled continuation (E0018)}
\textit{Continues an entry from an earlier page.}

Lavrus. cipally the fearcity of this plant which makes it valuable.

The benzoin tree will grow to a much larger size than the other, and its branches are more numerous. They are smooth, and of a fine light-green colour. The leaves on their upper surface are smooth and of a fine light-green colour, but their under surface is venose, and of a whitish cast. When bruised, they emit a fine fragrance. This tree was formerly mistaken for that which produces the drug called benzoin; which is now known to be obtained from a species of styrax. See Styrax.

The sassafras will grow to nearly the height of the others, though the branches are not so numerous. Its bark is smooth, and of a red colour, which beautifully distinguishes it in winter; whilst the fine shining green of its leaves constitutes its greatest beauty in summer. In these, indeed, there is a variety, and a very extraordinary one. Some are large, and of an oval figure; others are smaller, and of the same shape; whilst others again are so divided into three lobes, as to resemble the leaves of some sorts of the fig-tree. In America, the sassafras generally stands single in the woods, and along the fences round the fields. It flowers in May before the leaves come out; and being entirely covered with them, it is distinguished at a great distance by their beautiful yellow colour.

The root of the sassafras has a fragrant smell, and a sweetish, aromatic, subacrid taste; the bark tales much fronger than any other part, and the small twigs fronger than the large pieces. It is a warm aperient and corroborant, and frequently employed with good fuccels for purifying and sweetening the blood and juices. For these purpofes, infusions made from the rapped root or bark may be drank as tea. In some confiltutions indeed, such liquors are, by their fragrance, apt, on first taking them, to affect the head; but in such cafes they may be advantageously freed from their flavour by boiling. A decoction of fallafras, boiled down to the confluence of an extract, proves firmly bitterth and fubaltingent. Hoffman affures us, that he has frequently given this extract to the quantity of a temple a time, with remarkable facets, for strengthening the tone of the vificera in ca-cherts; is allo in the decline of intermittent fevers, and in hypochondriacal paffions. Saffras yields in di-filtration an extremely fragrant oil of a penetrating pungent taste, fo ponderous (notwithstanding lightened to the deep yield) as to find in water. Rec-tified spirit extracts the whole tate and fowl of falls-fres, and cleaves nothing in separation, hence the spiritous extract proves the most elegant and efficacious preparation, as containing the virtue of the root entire.

The bark of this tree is used by the women in Pennsylvania and other parts of North America in dyeing worsted a fine fading orange-colour, which does not fade in the sun. They use urine instead of alum in dyeing; and boil the dye in a brass boiler, because in an iron vessel it does not yield so fine a colour. The wood is made use of for polls belonging to the inclosures, for it is said to last a long time in the ground: but it is likewise said, that there is hardly any kind of wood which is more attacked by worms than this when it is exposed to the air without cover; and that in a short time it is quite worm-eaten through and through.

On cutting some part of the sassafras tree, or its shoots, and holding it to the nose, it has a strong but pleasant smell. Some people peel the root, and boil the peel with the beer which they are brewing, because they believe it wholesome. For the same reason, the peel is put into brandy either while it is distilling or after it is made. Professor Kalm informs us, that a decoction of the root of sassafras in water, drank every morning, is used with success in the dropsy.—When part of a wood is destined for cultivation, the sassafras trees are commonly left upon it, because they have a very thick foliage, and afford a cool shade to the cattle during the great heats. Some people get their bed-poles made of sassafras wood, in order to expel the bugs; for its strong scent, it is said, prevents those vermin from settling in them. For two or three years together this has the desired effect, or about as long as the wood keeps its strong aromatic smell; but after that time it has been observed to lose its effect. In Pennsylvania some people put chips of sassafras into their chests, where they keep all sorts of woollen stuffs, in order to expel the moths (or larvæ or caterpillars of moths or tinies) which commonly settle in them in summer. The root keeps its smell for a long while: Professor Kalm saw one which had lain five or six years in the drawer of a table, and still preferred the strength of its scent. The people also gather its flowers, and use it as tea.

The perfus, or alligator pear tree, is cultivated universally in the West Indies by all ranks of people. The fruit is pear-shaped, and from one to two pounds in weight. On removing a green skin or covering, we come to a yellow butyraceous substance; and in the heart find a large round feed or stone, which is unequal in the surface, and exceedingly hard and woody. This fruit is ripe in August and September, and constitutes one of the most agreeable articles of diet for six or eight weeks to the negroes. These pears, with a little fat and a plantain or two, afford a hearty meal. They are also served up at the tables of white people as choice fruit. When the pear is ripe, the yellow or eatable substance is firmer than butter, and tastes somewhat like butter or marrow: hence it is called by some the vegetable marrow. But however excellent this fruit is when ripe, it is very dangerous when pulled and eaten before maturity. Dr Wright says, he has repeatedly known it to produce fever and dysentery, which were removed with difficulty.—The leaves of this tree and the bark of the bead-vine or wild liquorice are made into pectoral decoctions by the common people.—The large stone is used for marking linen. The cloth is tied or held over the stone, and the letters are pricked out by a needle through the cloth and into the feed. The stain is a reddish brown, which never washes out.—The buds of the alligator tree are said to be used with succels as pifans against the venereal disease. An infusion of them in water, drank in the morning fasting, is strongly recommended for dislodging coagulated blood in the stomach produced by a fall or a few strokes on that important entrail. "The wild boars in the East Indies (says Labat) eat greedily of the mammees and avocado pears, which give their flesh a juicy and most agreeable savour."

a fulcious and most agreeable savour.
Cassia. The bark of this species is known in the
shops by the name of cassia lignea. This bark, which

Kalo's Travels in America.


\section*{Source record}
\begin{description}
\item[Scan ID] \texttt{britannica:309f765b1ecc83e0a7f7762e44b4a927a7555443:train:268}
\item[Image file] \texttt{page-0268-4f70fe5dafd6.jpg}
\item[Image SHA-256] \texttt{4f70fe5dafd60ccb4d258630071448612eebe999713dac3954c2011faacf3929}
\item[Dataset] NationalLibraryOfScotland/encyclopaedia\_britannica\_illustrated at revision \texttt{309f765b1ecc83e0a7f7762e44b4a927a7555443} (CC0-1.0)
\item[Extraction] model \texttt{nvidia/nemotron-3.5-parse}, raw reply \texttt{extract/raw/page-0268-4f70fe5dafd6.json}
\item[Raw reply SHA-256] \texttt{4ceda7d5b2e30d33bae6229226e4752869510e8902a370208fff59230ae24400}
\item[Correction] model \texttt{nvidia/Nemotron-3.5-Super-VL}, block crops, raw replies \texttt{correct/raw/crop\_nothink/page-0268-4f70fe5dafd6-b*.json}; super\_vl: 9
\item[Running header] LA U [ 6oz ] L AU
\end{description}
\end{document}
